% TOP_PROBLEM: {"problemId": "problem.superpolynomial-lower-bounds-for-frege-propositional-proof-systems", "canonicalTitle": "Superpolynomial Lower Bounds for Frege Propositional Proof Systems", "releaseRank": 104, "primaryDomain": "theoretical_computer_science", "primaryDomainLabel": "Theoretical computer science", "displayStatus": "open", "releaseStatus": "open"}
% !TeX program = lualatex
% Definition and sources only; this file is not an LLM solution attempt.
\documentclass[11pt]{article}
\usepackage[margin=25mm]{geometry}
\usepackage{fontspec}
\setmainfont{Latin Modern Roman}
\usepackage{amsmath,amssymb,mathtools}
\usepackage{unicode-math}
\setmathfont{Latin Modern Math}
\usepackage{xurl,hyperref}
\hypersetup{colorlinks=true,urlcolor=blue}
\setcounter{secnumdepth}{0}
\setlength{\parindent}{0pt}
\setlength{\parskip}{5pt}
\setlength{\emergencystretch}{3em}
\begin{document}
\section*{\texorpdfstring{Superpolynomial Lower Bounds for Frege Propositional Proof Systems}{Superpolynomial Lower Bounds for Frege Propositional Proof Systems}}\label{104}
\subsection{Definitions and mathematical statement}
A Frege system is a fixed sound and complete finite propositional axiom-and-rule calculus with substitution instances; proof size counts all written symbols. For a tautology $\tau$, let $s_F(\tau)$ be the minimum size of an $F$-proof. Seek an explicit growing family $\tau_n$, conventionally efficiently generated from $n$, for which
\[
 s_F(\tau_n)\text{ is not bounded by any polynomial in }|\tau_n|.
\]
A corresponding line-count lower bound is another allowed target in the catalogue. The family must consist of tautologies, and the system is a standard unrestricted Frege system, not merely a depth-restricted or otherwise weakened proof calculus.
\par\smallskip\textit{Formulation and concept references:} \href{https://doi.org/10.1017/9781108643818}{[S1]}.

\subsection{Short English statement}
Can one explicitly exhibit tautologies requiring superpolynomially long proofs in ordinary Frege reasoning systems?
\subsection{Sources}
\begin{itemize}
\item[S1]J. Krajicek, Proof Complexity, Cambridge University Press, 2019.
\url{https://doi.org/10.1017/9781108643818}.
\textit{Formulation source.}
\end{itemize}
\end{document}
